\section{Evaluation am Fallbeispiel} \label{evaluation}

Dieses Kapitel demonstriert die Pipeline am konkreten Migrationsszenario und beschreibt, wie die Vollständigkeit der erzeugten Artefakte sichergestellt wird.

\subsection{Auswahl des Szenarios und Quell-Moduls}

Als Demonstrationsszenario dient die Migration eines Node.js-basierten Lizenz-Checkers nach Go. Das Modul wurde gewählt, weil es ausreichend fachliche und technische Komplexität besitzt, um einen vollständigen Durchstich durch alle Pipelinephasen zu zeigen, ohne den Rahmen der Projektarbeit zu überschreiten. Es umfasst eine Kommandozeilenschnittstelle, Datei- und Verzeichnisverarbeitung, die Auswertung von Paketmetadaten sowie die Ausgabe strukturierter Lizenzberichte. Diese Funktionalitäten lassen sich klar in einzelne Use Cases zergliedern, sodass die Spec-Writer-Fan-out-Logik des Orchestrators gezielt demonstriert werden kann. Das Modul enthält externe Abhängigkeiten, die ein konkretes Bibliotheks-Mapping zwischen Node.js- und Go-Ökosystem erfordern, und besitzt hinreichend definiertes Ein- und Ausgabeverhalten für einen manuellen Input-Output-Vergleich im Evaluationsschritt.

\subsection{Durchführung des spezifikationsgetriebenen Workflows}

Dieser Abschnitt dokumentiert den Ausführungsstand der Pipeline. Die drei Phasen wurden vollständig durchlaufen und alle Artefakte erzeugt. Der abschließende Black-Box-Vergleich zwischen Quell- und Zielsystem wurde manuell nach dem in Abschnitt~\ref{sec:funktionale-validierung} beschriebenen Verfahren durchgeführt; die Befunde sind in Abschnitt~\ref{sec:io-befunde} dokumentiert. Im Zuge des Vergleichs traten zwei methodisch eigenständige Defektklassen zutage, die in iterativen Skill-Anpassungen aufgearbeitet wurden (Abschnitte~\ref{sec:builder-gap} und~\ref{sec:datafluss-gap}).

Im Fallbeispiel wurde die vollständige Pipeline ausgeführt. Der Orchestrator steuerte die Übergaben zwischen den sechs fachlichen Subagenten (Analyzer, Spec-Writer, Architect, Planner, Task Decomposer, Builder) sowie dem Verifier als unabhängiger Prüfinstanz, sodass Analyse, Spezifikation, Architektur, Planung, Aufgabenzerlegung und Code-Synthese vollständig durchgeführt wurden. Der abschließende I/O-Vergleich gegen das Quellsystem wurde anschließend manuell durchgeführt; die Befunde sind in Abschnitt~\ref{sec:io-befunde} dokumentiert.

In Phase~1 erschließt der Analyzer die Quellbasis und erzeugt \texttt{ANALYSIS.md} und \texttt{FUNCTIONALITY-INDEX.md}. Die ausgeführte Pipeline identifizierte acht Funktionalitätselemente: F-001 (CLI-Optionsmodell und Normalisierung), F-002 (Abhängigkeitsgraph-Laden und Traversierungssteuerung), F-003 (Abflachen des Graphen zu einem Modulinventar), F-004 (Lizenzerkennungs- und Normalisierungspipeline), F-005 (Filterung, Paketeinschränkung und Policy-Exits), F-006 (Ausgabe-Rendering und Persistenz), F-007 (Custom-Format und JSON-Konfigurationsladung) sowie F-008 (Stack-Utility-Lifecycle). Alle acht wurden vom Analyzer als \emph{ready} klassifiziert. Der Orchestrator löste daraufhin acht parallele Spec-Writer-Läufe aus, einer pro Funktionalitätselement, und sammelte die Ergebnisse in \texttt{SPEC-INDEX.md}.

In Phase~2 leitete der Architect aus der konsolidierten Spezifikationsbasis eine Go-Zielarchitektur ab: Paketstruktur, Fehlerbehandlungskonventionen, geeignete Standardbibliotheksfunktionen und externe Go-Bibliotheken wurden mit Context7-Unterstützung festgelegt. Der Planner überführte diese Entscheidungen anschließend in einen geordneten Migrationsplan. In Phase~3 zerlegte der Task Decomposer den Plan in umsetzbare Arbeitseinheiten und der Builder erzeugte Go-Zielcode, der in \texttt{BUILD.md} mit Verifikationsnachweis dokumentiert ist. Alle acht Funktionalitäten wurden vollständig durch die Pipeline geführt; der abschließende Vergleich des beobachtbaren CLI-Verhaltens zwischen Quell- und Zielsystem erfolgte anschließend manuell nach dem in Abschnitt~\ref{sec:funktionale-validierung} beschriebenen Verifikationsplan; die Befunde sind in Abschnitt~\ref{sec:io-befunde} dokumentiert.

\subsection{Beispiel: Spezifikation und Zielcode im Vergleich}

Die folgenden Listings zeigen den tatsächlichen Output der Pipeline für die Funktionalität \textbf{F-004: License detection and normalization pipeline}. Listing~\ref{lst:spec} gibt einen Auszug aus der vom Spec-Writer erzeugten \texttt{SPEC.md}, Listing~\ref{lst:target} den daraus vom Builder generierten Go-Zielcode.

\begin{lstlisting}[style=codestyle, caption={Auszug aus der erzeugten SPEC.md fuer F-004 (License detection and normalization pipeline)}, label=lst:spec]
## Requirements

- FR-001: System MUST resolve license from metadata fields
          (`license`, `licenses`) before fallback mechanisms.
- FR-002: System MUST apply fallback precedence:
          metadata -> README-derived -> license-file-derived.
- FR-004: System MUST mark heuristically inferred values
          with trailing `*` to distinguish inferred vs direct.
- FR-005: System MUST normalize URL/file-style references
          to `Custom: <value>`.
- FR-006: File fallback MUST follow deterministic order:
          LICENSE*, LICENCE*, COPYING, README.

## Acceptance Scenario (P1)

Given a module with a valid SPDX `license` string,
When the pipeline resolves license data,
Then that SPDX value is returned without file-based fallback.

## Open Decisions

- OD-001: When `license` and `licenses` conflict, confirm
  exact legacy tie-break order from source code/tests.
- OD-002: Confirm exact fallback value when no license
  signal is found across metadata, README, and files.
\end{lstlisting}

\begin{lstlisting}[style=codestyle, language=Go, caption={Vom Builder erzeugter Go-Zielcode: license/detect.go}, label=lst:target]
package license

import "theseus/target/internal/core"

type Detector struct{}

func (Detector) Resolve(
    input core.ModuleLicenseInput,
) (core.NormalizedLicenseValue, string) {

    // FR-001/FR-002: metadata first
    if v := resolveFromMetadata(input); v != "" {
        return core.NormalizedLicenseValue(v), ""
    }
    // FR-002: README fallback
    if v, _ := normalize(input.ReadmeText); v != "" {
        return core.NormalizedLicenseValue(v), "README"
    }
    // FR-002/FR-006: deterministic file fallback
    for _, candidate := range orderedCandidates(
        input.CandidateFiles,
    ) {
        if v, _ := normalize(candidate.Content); v != "" {
            return core.NormalizedLicenseValue(v),
                   candidate.Name
        }
    }
    // OD-002: sentinel until legacy parity confirmed
    return core.NormalizedLicenseValue("UNKNOWN"), ""
}
\end{lstlisting}

Die Listings zeigen den tatsächlichen Output der ausgeführten Pipeline. Die Anforderungen FR-001 bis FR-006 aus \texttt{SPEC.md} sind direkt auf Implementierungsschritte im Go-Zielcode abbildbar: Die Fallback-Reihenfolge (Metadata $\to$ README $\to$ Dateien) entspricht den in der Spezifikation nummerierten Teilschritten, und die beiden offenen Entscheidungen OD-001 und OD-002 sind als Kommentare im generierten Code explizit markiert. Das zeigt, wie der spec-getriebene Ansatz Unsicherheiten aus der Analyse nicht wegabstrahiert, sondern als \emph{NEEDS CLARIFICATION}-Marker durch die gesamte Pipeline trägt, bis sie manuell oder durch weitere Analyse aufgelöst werden können.

\subsection{Funktionale Validierung und Feature Parity} \label{sec:funktionale-validierung}

Die funktionale Validierung folgt einem Black-Box-Ansatz im Sinne des klassischen \emph{Differential Testing}, bei dem zwei Implementierungen desselben Programms mit identischen Eingaben ausgeführt und ihre Ausgaben verglichen werden, ohne dass Annahmen über die interne Struktur erforderlich sind \parencite{mckeeman1998differentialtesting}. Für die LLM-gestützte Code-Übersetzung hat dieses Vorgehen unter dem Begriff der I/O-Äquivalenz aktuelle Aufmerksamkeit erfahren \parencite{eniser2025realworld}. In der vorliegenden Arbeit werden Quell- und Zielimplementierung als geschlossene CLI-Werkzeuge betrachtet und an identischen Eingaben gegeneinander ausgeführt; die Beurteilung der Feature Parity ergibt sich aus dem Vergleich der beobachtbaren Ausgaben. Als Quellsystem dient \texttt{davglass/license-checker}, dessen sechs Use Cases von der Auflistung aller Dependency-Lizenzen über die Erzeugung maschinenlesbarer Compliance-Artefakte bis hin zur Policy-Durchsetzung in CI-Umgebungen reichen und in den Spezifikationen vollständig abgebildet sind. Statische Analysen und Tests des Zielsystems (\texttt{go build}, \texttt{go vet}, \texttt{go test}) dienen lediglich als Vorab-Filter, der sicherstellt, dass das Zielbinär ausführbar und gegen die Akzeptanzszenarien der Spezifikationen abgesichert ist; die Aussage über die Verhaltensgleichheit gegenüber dem Quellsystem entsteht erst durch den nachgelagerten manuellen I/O-Vergleich. Durchführung und Befunde sind in Abschnitt~\ref{sec:io-befunde} dokumentiert.

\subsubsection{Verifikationsverfahren pro Funktionalität}

Damit Abweichungen sich nicht in einem Gesamtoutput verschleiern, wird der Vergleich getrennt nach den acht identifizierten Funktionalitäten geführt. Wo das Zielsystem während Phase~3 bereits Integrationstests mit Golden-Files erzeugt hat (etwa \texttt{render\_export} oder \texttt{license\_resolution\_compat}), werden diese Fixtures wiederverwendet und gegen die Outputs des Quellsystems gespiegelt, statt die Vergleichsmenge neu aufzubauen. Tabelle~\ref{tab:io-vergleich} fasst das Verifikationsverfahren und die konkreten Testszenarien pro Funktionalität zusammen.

\begin{small}
\begin{longtable}{|p{0.7cm}|p{4.0cm}|p{10.5cm}|}
\caption{Black-Box-Verifikationsverfahren je Funktionalität für den I/O-Vergleich zwischen \texttt{davglass/license-checker} (Node.js) und der Go-Zielimplementierung. Wo vorhandene Integrationstests und Golden-Files des Zielsystems verwendbar sind, werden sie als Vergleichsbasis wiederverwendet.} \label{tab:io-vergleich} \\
\hline
\textbf{F} & \textbf{Funktionalität} & \textbf{Black-Box-Verifikationsverfahren} \\
\hline
\endfirsthead

\multicolumn{3}{l}{\textit{Fortsetzung von Tabelle~\ref{tab:io-vergleich}}} \\
\hline
\textbf{F} & \textbf{Funktionalität} & \textbf{Black-Box-Verifikationsverfahren} \\
\hline
\endhead

\hline
\multicolumn{3}{r}{\textit{Fortsetzung auf der nächsten Seite}} \\
\endfoot

\hline
\endlastfoot

001 & CLI-Optionsmodell & Vergleich von Stdout, Stderr und Exit-Code für alle dokumentierten Flag-Kombinationen, für \texttt{-{}-help}/\texttt{-{}-version} sowie für gezielt ungültige Flag-Kombinationen. \\
\hline
002 & Abhängigkeitsgraph laden & Drei Referenz-Repositories durchlaufen: ein flaches Projekt mit wenigen Direkt-Abhängigkeiten, ein mittelgroßes Projekt mit Dev-/Peer-Abhängigkeiten und ein Projekt mit zyklischen Verweisen. Vergleich der traversierten Modulmengen. \\
\hline
003 & Graph abflachen & Vergleich des sortierten Modulinventars über den bereits vorhandenen \texttt{flatten\_contract}-Integrationstest des Zielsystems. Deduplizierung und Zirkel-Guard werden zusätzlich über konstruierte Sonderfälle geprüft. \\
\hline
004 & Lizenzerkennung & Wiederverwendung der \texttt{license\_resolution\_compat}- und \texttt{license\_resolution\_fallback}-Fixtures: Testpakete mit nur \texttt{license}-Metadatum, nur README-Hinweis, nur LICENSE-Datei, allen drei kombiniert sowie ohne Signatur. Verifiziert werden Precedence-Reihenfolge, \texttt{*}-Marker und \texttt{Custom:}-Normalisierung gegen die Outputs des Quellsystems. \\
\hline
005 & Filter und Policy & Wiederverwendung der Fixtures aus \texttt{filter\_policy\_integration}: gezielte Konfigurationen für \texttt{-{}-failOn}, \texttt{-{}-onlyAllow} und Private-Filter (Manifest mit \texttt{private: true}). Verifiziert werden Exit-Codes, Filtermengen und Policy-Verhalten bei verletzten Bedingungen. \\
\hline
006 & Ausgabe-Rendering & Für ein Referenz-Repository wird jedes Ausgabeformat erzeugt und gegen die im Zielsystem hinterlegten Golden-Files (\texttt{json.golden}, \texttt{csv.golden}, \texttt{markdown.golden}, \texttt{tree.golden}, \texttt{summary.golden}) sowie spiegelbildlich gegen die Quellausgabe verglichen; reine Formatierungsunterschiede werden als semantisch identisch klassifiziert. \\
\hline
007 & Custom-Format & Wiederverwendung der Fixtures aus \texttt{custom\_format\_inline} und \texttt{custom\_format\_exclusion}: identische Custom-JSON-Konfiguration in beide Systeme laden; Vergleich der gerenderten Felder und der Default-Wert-Auflösung bei unvollständiger Konfiguration. \\
\hline
008 & Stack-Utility-Lifecycle & Interner Hilfsdatentyp ohne direkt beobachtbare CLI-Schnittstelle; indirekt abgedeckt durch korrektes Verhalten der nachgelagerten Funktionalitäten F-002 und F-003. \\

\end{longtable}
\end{small}

\subsubsection{Befunddokumentation und Fehlerzuordnung}

Pro Funktionalität wird ein Vergleichsbefund festgehalten: \emph{identisch}, \emph{semantisch identisch} bei reinen Formatabweichungen, oder \emph{abweichend} mit konkreter Beschreibung. Tritt eine Abweichung auf, wird sie auf die ursprüngliche Pipelinephase zurückgeführt. Liegt der Befund daran, dass die zugrunde liegende \texttt{SPEC.md} eine relevante Anforderung nicht eindeutig festgelegt hat, gilt der Befund als \emph{Spec-Lücke}; war die Anforderung dort eindeutig formuliert, der Zielcode weicht aber dennoch ab, gilt er als \emph{Builder-Fehler}. Diese Zuordnung ist diagnostisch wichtig, weil sie zeigt, an welcher Stelle der Pipeline (Spezifikations- oder Implementierungsphase) eine Korrektur ansetzen muss, und damit die Aussagekraft der Verifikation über reine Funktionsergebnisse hinaus erhöht.

Der Vergleich wird im PoC-Rahmen bewusst manuell geführt, weil die acht Funktionalitäten überschaubar sind und der diagnostische Wert in der Phasenzuordnung liegt, nicht in der Skalierung der Vergleichsmenge. Eine automatisierte Variante über einen Cross-Language-Differenzfuzzer wie Fluorine~\parencite{eniser2025realworld} ist für die Bachelorarbeit als Erweiterung vorgesehen; sie würde den hier definierten Verifikationsplan in einen wiederholbaren Prüfrahmen überführen und denselben Befund-Schema-Aufbau übernehmen.

\subsubsection{Durchführung und Befunde des manuellen Vergleichs} \label{sec:io-befunde}

Der Vergleich wurde nach einem festen Schema pro Funktionalität durchgeführt. Für jeden Vergleichsfall wird ein minimales Node.js-Eingabeprojekt mit den für die jeweilige Funktionalität relevanten Eigenschaften vorbereitet (\texttt{package.json} mit definierten Abhängigkeiten und Lizenzfeldern, einmaliges \texttt{npm install}). Anschließend wird das Quellsystem \texttt{davglass/license-checker} über \texttt{npx -{}-yes license-checker} mit den vereinbarten Argumenten im Eingabeverzeichnis ausgeführt und \texttt{stdout}, \texttt{stderr} sowie Exit-Code festgehalten. Das aus \texttt{target/cmd/licensechecker} kompilierte Go-Binary wird daraufhin mit denselben Argumenten im selben Verzeichnis ausgeführt und die Ergebnisse werden analog protokolliert. Die beobachtbaren Ausgaben werden gegeneinander gestellt und einer der drei Äquivalenzklassen aus Abschnitt~\ref{sec:funktionale-validierung} zugeordnet; bei Abweichungen erfolgt die Klassifikation als \emph{Spec-Lücke} oder \emph{Builder-Fehler}.

Für den Durchstich wurden zehn Vergleichsfälle vorbereitet, die alle acht Funktionalitätselemente außer F-008 (interner Hilfsdatentyp ohne CLI-Oberfläche) adressieren. Die hier dargestellten Befunde beziehen sich auf den Pipelinedurchlauf nach der in Abschnitt~\ref{sec:builder-gap} dokumentierten ersten Skill-Iteration; die ursprüngliche Pipelineversion vor dieser Anpassung erzeugte kein lauffähiges Binary und konnte daher nicht direkt verglichen werden. Tabelle~\ref{tab:io-befunde} fasst die Befunde des durchgeführten Vergleichs zusammen.

\begin{table}[H]
\centering
\begin{small}
\begin{tabular}{|p{1.3cm}|p{3.5cm}|p{2.6cm}|p{5.5cm}|}
\hline
\textbf{F-ID} & \textbf{Vergleichsfall} & \textbf{Befund} & \textbf{Anmerkung} \\
\hline
F-001 & Versionsausgabe (\texttt{-{}-version}) & Diff vorhanden & Erwartet; Versionsangaben unterscheiden sich naturgemäß. \\
\hline
F-001 & Hilfetext (\texttt{-{}-help}) & Diff vorhanden & Erwartet; das Hilfetext-Layout der \texttt{cobra}-basierten Go-Implementierung weicht vom \texttt{commander.js}-Layout des Quellsystems ab. \\
\hline
F-002 & Flaches Projekt (\texttt{-{}-json}) & Semantisch identisch & Beide Systeme listen Root-Paket und direkte Abhängigkeit mit übereinstimmenden Lizenzangaben. \\
\hline
F-003 & Transitive Deduplizierung (\texttt{-{}-json}) & Semantisch identisch & Modulmenge und Lizenzangaben stimmen über die gesamte Abhängigkeitskette überein. \\
\hline
F-004 & Minimales Projekt mit \texttt{-{}-json} & Semantisch identisch & JSON-Top-Level-Schema und Schlüsselbenennung übereinstimmend. \\
\hline
F-005 & \texttt{-{}-failOn GPL} ohne GPL-Dep & Semantisch identisch & Beide Systeme listen die Abhängigkeitskette ohne Policy-Verletzung mit Exit-Code $0$. \\
\hline
F-005 & \texttt{-{}-onlyAllow MIT} & Abweichend (Spec-Lücke) & Quellsystem bricht ab, weil es \texttt{"private": true} unbedingt als \texttt{UNLICENSED} interpretiert; das Zielsystem liest das \texttt{license}-Feld wörtlich. Die zugrundeliegende Regel des Quellsystems ist in der \texttt{SPEC.md} nicht als unbedingter Override abgebildet (siehe Abschnitt~\ref{sec:datafluss-gap}). \\
\hline
F-006 & CSV-Ausgabe & Semantisch identisch & Spalten-Header und Datenzeilen übereinstimmend. \\
\hline
F-006 & Markdown-Ausgabe & Semantisch identisch & Layout und Modulauflistung übereinstimmend. \\
\hline
F-007 & Custom-Format (\texttt{-{}-customPath}) & Semantisch identisch & \texttt{custom-format.json} wird vom Zielsystem geladen und in die Render-Pipeline propagiert. \\
\hline
F-008 & -- & n.\,a. & Interner Hilfsdatentyp ohne CLI-Oberfläche; indirekt über F-002 und F-003 abgedeckt. \\
\hline
\end{tabular}
\end{small}
\caption{Befunde des manuellen I/O-Vergleichs zwischen \texttt{davglass/license-checker} und der Go-Zielimplementierung nach der ersten Skill-Iteration}
\label{tab:io-befunde}
\end{table}

Die Befunde lassen drei methodisch belastbare Aussagen zu. Erstens verhält sich die CLI-Oberfläche des Zielsystems erwartungsgemäß: \texttt{-{}-version} und \texttt{-{}-help} liefern jeweils eine Ausgabe und einen sinnvollen Exit-Code; die Divergenzen gegenüber dem Quellsystem sind auf der Ebene des Werkzeugrahmens (\texttt{cobra} statt \texttt{commander.js}) zu erklären und nicht funktional. Zweitens decken acht von neun inhaltlichen Vergleichsfällen das Verhalten des Quellsystems semantisch ab -- über alle Pipeline-durchschneidenden Funktionalitäten hinweg (F-002 bis F-007), inklusive der zuvor defekten Loader-zu-Flatten-Verkabelung und der zuvor abweichenden Render-Schemata. Der nach der ersten Skill-Iteration verbleibende Befund auf F-005 \texttt{-{}-onlyAllow MIT} ist von den restlichen abweichenden Fällen strukturell verschieden: er ist weder Verkabelungs- noch Renderdefekt, sondern eine im Quellcode des Quellsystems unbedingt wirkende Werttransformation (\texttt{private: true} setzt das Feld \texttt{licenses} unabhängig vom CLI-Flag auf \texttt{UNLICENSED}), die in der flagzentrierten Rekonstruktion der ursprünglichen Pipelinedurchläufe nicht erfasst wurde und in Abschnitt~\ref{sec:datafluss-gap} als zweite methodisch eigenständige Defektklasse aufgearbeitet wird. Drittens ist methodisch bemerkenswert, dass keiner der ursprünglichen Defekte durch die im Builder-Output enthaltenen Integrationstests erfasst wurde -- alle Tests waren grün, ebenso \texttt{go build ./...} und \texttt{go vet ./...}. Die Defekte zeigten sich erst, wenn das End-to-End-Verhalten gegen das Quellsystem geprüft wurde. Damit bestätigt der manuelle I/O-Vergleich die in Kapitel~\ref{grundlagen} angelegte These, dass funktionale Korrektheit eines aus Spezifikationen generierten Systems erst auf der Ebene beobachtbarer Ausgaben verlässlich beurteilt werden kann, und nicht aus dem Bestehen lokaler Test- und Build-Schritte abgeleitet werden darf.

\subsubsection{Erste Skill-Iteration: Verkabelungslücke des Builder-Outputs} \label{sec:builder-gap}

Im ursprünglichen Pipelinedurchlauf trat ein methodisch relevanter Befund auf, der den in Abschnitt~\ref{sec:agenten} eingeführten Verifier-Mechanismus direkt betrifft: Der Builder-Agent hatte für die acht Funktionalitäten Library-Code unter \texttt{target/internal/} sowie Integrationstests unter \texttt{target/test/integration/} erzeugt, jedoch keinen lauffähigen Entry-Point in Form eines \texttt{main.go} unter \texttt{target/cmd/}. Damit fehlte die End-to-End-Verkabelung der Stufen \emph{Loader}~$\rightarrow$~\emph{Flatten}~$\rightarrow$~\emph{License-Resolution}~$\rightarrow$~\emph{Filter}~$\rightarrow$~\emph{Render}, ohne die ein I/O-Vergleich gegen das Quellsystem nicht möglich gewesen wäre. Ein zur Überbrückung manuell ergänztes \texttt{main.go} bestätigte den Befund empirisch: Alle Vergleichsfälle, die einen vollständigen Pipeline-Durchstich erfordern, erkannten ausschließlich das Root-Modul; zusätzlich produzierten die Render-Stufen für \texttt{-{}-json}, \texttt{-{}-csv} und Markdown jeweils ein vom Quellsystem abweichendes Schema, und \texttt{-{}-customPath} wurde vollständig ignoriert.

Auffällig ist, dass dieser Mangel vom Verifier in der ursprünglichen Konfiguration nicht erkannt wurde: \texttt{go build ./...} und \texttt{go vet ./...} laufen auch ohne Entry-Point fehlerfrei durch, weil die Werkzeuge nur die vorhandenen Pakete kompilieren. Die Verifier-Checkliste enthielt keinen expliziten Check auf ein lauffähiges Binary für CLI-orientierte Use Cases. Dies passt zu den in Abschnitt~\ref{sec:agenten} angeführten Befunden zur Selbstüberschätzung agentischer Systeme \parencite{huang2024cannotselfcorrect,korn2026reportingprompting}: Eine Phase galt artefaktbasiert als bestanden, obwohl ein für die Akzeptanzszenarien zentrales Artefakt fehlte. Der Befund stützt damit empirisch das Argument, warum die Verifier-Rolle vom Orchestrator getrennt und mit einer expliziten, granularen Checkliste ausgestattet sein muss.

Als unmittelbare Konsequenz wurden die Pipeline-Skills entlang der natürlichen Pipeline-Flussrichtung erweitert, statt einzelne Agentenprofile mit Spezialfällen zu beschweren. Im \texttt{cli-analyzer}-Skill wurde verankert, dass ANALYSIS.md für jedes Output-Format des Quellsystems eine reale, durch tatsächliche Ausführung gewonnene Beispielausgabe (\emph{Golden Output Snippet}) sowie eine explizite \texttt{End-to-End-Pipeline}-Sektion enthalten muss; das End-to-End-Verhalten ist damit aber bewusst keine eigene Funktionalität, sondern eine Architekturanforderung. Im \texttt{go-architect}-Skill wurde im Gegenzug die Pflicht eines verbindlichen \texttt{Composition-Root}-Abschnitts in ARCHITECTURE.md verankert, der Entry-Point-Pfad, Stage-Sequenz, konkrete Stage-Konstruktoren, eine Flag-zu-Stage-Propagationstabelle sowie eine Akzeptanzreferenz auf die Golden Output Snippets der ANALYSIS.md festschreibt. Damit hat der Builder bereits aus den Architektur-Artefakten eine eindeutige Bauanleitung für \texttt{cmd/<tool>/main.go}. Im \texttt{orchestrator-completion}-Skill wurde die Phasen-Checkliste um genau diese beiden Punkte ergänzt: Phase~2 prüft das Vorhandensein des Composition Roots, Phase~3 prüft, ob das Binary baut, auf \texttt{-{}-help} und \texttt{-{}-version} reagiert und ob seine End-to-End-Ausgabe den Golden Output Snippets semantisch entspricht. Das Verifier-Profil bleibt generisch und verweist auf den Skill als verbindliche Quelle der Checkliste. Ergänzend wurde im Builder-Profil festgeschrieben, dass für CLI-orientierte Use Cases ein \texttt{cmd/<tool>/main.go} als Pflichtartefakt zur BUILD.md gehört. Alle Anpassungen wurden in den jeweiligen Skill- und Profildateien unter \texttt{.agents/skills/} und \texttt{.opencode/agents/} festgeschrieben.

Ein anschließender erneuter Pipelinedurchlauf mit den erweiterten Skills bestätigt die Wirksamkeit der Anpassungen empirisch. Der Builder erzeugte ohne externe Ergänzung ein lauffähiges Binary unter \texttt{target/cmd/licensechecker/main.go}, das die im Composition Root spezifizierte Stage-Sequenz instanziiert und die in der Flag-zu-Stage-Propagationstabelle festgehaltenen Konfigurationen propagiert. Die im selben Lauf produzierten I/O-Vergleichsergebnisse (Tabelle~\ref{tab:io-befunde}) zeigen, dass die zuvor identifizierten Verkabelungs- und Schemadefekte (F-002, F-003, F-004, F-005 \texttt{-{}-failOn}, F-006 CSV, F-006 Markdown, F-007) ohne Ausnahme behoben sind. Die Skill-Erweiterung hat damit nicht eine einzelne Diagnose adressiert, sondern eine ganze Defektklasse strukturell geschlossen.

\subsubsection{Zweite Skill-Iteration: Daten-Fluss-Lücke der Quellrekonstruktion} \label{sec:datafluss-gap}

Der nach der ersten Iteration verbleibende Befund auf F-005 \texttt{-{}-onlyAllow MIT} weist auf eine zweite, methodisch eigenständige Defektklasse hin. Das Quellsystem bricht mit Exit-Code $1$ und der Meldung \texttt{"... is licensed under \"UNLICENSED\" which is not permitted by the -{}-onlyAllow flag"} ab; das Zielsystem hingegen verarbeitet die Eingabe ohne Policy-Verletzung weiter, weil es das deklarierte Lizenzfeld (\texttt{"license": "MIT"}) wörtlich übernimmt. Die Ursache liegt im Quellsystem in einer einzigen unbedingt wirkenden Werttransformation: In \texttt{lib/index.js} wird der Wert von \texttt{licenses} bei Paketen mit \texttt{"private": true} unabhängig vom CLI-Flag und unabhängig vom deklarierten Lizenzwert auf \texttt{UNLICENSED} gesetzt. Der dazugehörige Test (\texttt{tests/test.js}, \emph{should recognise private modules}) prüft genau dieses Verhalten an einer Fixture mit \texttt{"private": true} ohne \texttt{license}-Feld.

In der durch den Pipelinedurchlauf erzeugten Spezifikation taucht der Marker \texttt{private} mehrfach auf, jedoch ausschließlich im Kontext des Flags \texttt{-{}-excludePrivatePackages}: ANALYSIS.md führt ihn als Filter-Konzept; die SPEC.md des Filter-Slices formuliert die zugehörige FR als \emph{System MUST remove private packages when excludePrivatePackages is enabled}. Die unbedingte Werttransformation des Felds \texttt{licenses} -- die ohne jedes Flag wirkt -- ist in keiner SPEC.md erfasst. Strukturell handelt es sich also nicht um einen Rekonstruktionsfehler im engeren Sinne, sondern um eine systematische Blindstelle der bisher verwendeten Rekonstruktionsachse: Die \texttt{cli-analyzer}-Skill in ihrer ersten Fassung schreibt eine Rekonstruktion \emph{entlang der CLI-Oberfläche} vor (Flag~$\rightarrow$~Verhalten~$\rightarrow$~Test). Eine flagunabhängige Werttransformation auf einem Eingabefeld besitzt diese Achse nicht und bleibt damit unsichtbar; der Test, der das Verhalten exerziert, wird implizit demjenigen Slice zugeschlagen, in dessen Flagname das Feld vorkommt.

Diese zweite Defektklasse ist abstrakt adressierbar, indem die Rekonstruktion um eine \emph{datenflussorientierte} Achse ergänzt wird. Im \texttt{cli-analyzer}-Skill wurde dafür eine zusätzliche Pflichtsektion \texttt{\#\# Input-Field Inventory} verankert. Diese verlangt für jedes Eingabefeld, das das Quellsystem aus externen Quellen (Konfigurationsdateien, Manifest-Dateien, Paket-Deskriptoren, Umgebungsvariablen, Stdin, Netzwerkantworten) liest, eine explizite Auflistung: (a) Feldname und Quelle, (b) Defaultpfad durch die Pipeline, (c) jeder konditionale Override mit Code-Quelle und auslösender Bedingung -- insbesondere ausgewiesen, ob er von einem CLI-Flag, einer Umgebungsvariablen oder einem Konfigurationsschlüssel gesteuert wird oder unbedingt wirkt, und (d) eine Querreferenz zu Tests im Quellsystem, die den Override-Pfad exerzieren (auch dann, wenn deren \texttt{describe}-Block dem Namen nach zu einem anderen Slice zu gehören scheint). Für jeden unbedingten Override muss der Analyzer zusätzlich eine Probe-Fixture konstruieren, das Quellsystem darauf ausführen und das Ergebnis als zusätzliches \emph{Golden Output Snippet} hinterlegen. Im \texttt{orchestrator-completion}-Skill prüft Phase~1 das Vorhandensein und die Vollständigkeit dieses Inventars sowie die Forderung, dass jeder unbedingte Override als eigene FR in der SPEC.md desjenigen Slices erscheint, der das betroffene Ausgabefeld besitzt -- und nicht nur in dem Slice, der zufällig denselben Namensbestandteil im CLI-Flag trägt.

Methodisch ordnet sich diese zweite Iteration als orthogonale Ergänzung zur ersten ein: Die erste Iteration adressierte die \emph{Control-Flow-Lücke} zwischen Spezifikation und lauffähigem Binary; die zweite Iteration adressiert die \emph{Data-Flow-Lücke} zwischen Quellverhalten und Spezifikation. Beide Defektklassen waren durch das Bestehen lokaler Test- und Build-Schritte nicht zu erkennen, beide wurden erst durch das beobachtungsbasierte Differential Testing aufgedeckt, und beide ließen sich strukturell durch Anpassungen an den Pipeline-Skills schließen, ohne dass eine domänenspezifische Sonderbehandlung in den Agentenprofilen verankert werden musste. Der verbliebene FAIL ist damit nicht das Ergebnis eines unzureichenden Builder-Outputs, sondern der empirische Beleg dafür, dass eine ausschließlich kontrollflussorientierte Reverse-Engineering-Methode eine systematische und benennbare Klasse von Verhaltensspezifikationen übersieht; eine empirische Verifikation der Wirksamkeit der zweiten Iteration analog zur ersten würde einen erneuten Pipelinedurchlauf mit den weiter erweiterten Skills erfordern und ist außerhalb des zeitlichen Rahmens dieser Projektarbeit verortet.

\subsection{Qualitätsdimensionen der erzeugten Artefakte}

Die Evaluation des PoC erstreckt sich auf zwei Artefaktklassen: die generierten Spezifikationen und den erzeugten Go-Zielcode. Für beide werden im Folgenden die maßgeblichen Qualitätsdimensionen beschrieben, die im Evaluationsdurchlauf beurteilt werden.

\subsubsection{Qualitätsdimensionen der Spezifikationsartefakte}

Wie \textcite{midolo2026promptguidelines} zeigen, sind für hochwertige Spezifikationen vor allem explizite I/O-Angaben, Vor- und Nachbedingungen sowie Beispiele relevant. Für die vorliegende Evaluation werden folgende Dimensionen herangezogen:

\textbf{Präzision:} Sind I/O-Formate, Datentypen, Vor- und Nachbedingungen sowie Akzeptanzszenarien so konkret formuliert, dass der Builder ohne Rückfragen implementieren kann? Vage oder implizite Anforderungen gelten als Qualitätsmangel.

\textbf{Rückverfolgbarkeit:} Verweist jede funktionale Anforderung auf eine konkrete Evidenzstelle in \texttt{ANALYSIS.md}? Lücken in der Rückverfolgbarkeit erschweren die Fehlersuche bei Abweichungen zwischen Quell- und Zielverhalten.

\textbf{Vollständigkeit:} Deckt die Spezifikation alle in \texttt{FUNCTIONALITY-INDEX.md} identifizierten Funktionalitätselemente ab? Fehlende Slices bedeuten blinde Flecken in der Neuimplementierung.

\textbf{Offene-Fragen-Behandlung:} Werden unklare Punkte explizit als offene Fragen markiert statt still angenommen? Implizite Annahmen gelten als Qualitätsmangel im Sinne des Spec-Quality-Skills \parencite{feng2024abstain}.

Für die Aussagekraft der Evaluation ist zudem relevant, wie transparent Prompt- und Kontextentscheidungen dokumentiert werden, da \textcite{korn2026reportingprompting} deutliche Reproduzierbarkeitslücken in der SE-Forschung aufzeigen.

\subsubsection{Qualitätsdimensionen des Zielcodes}

\textbf{Korrektheit (I/O-Äquivalenz):} Erzielt der generierte Go-Code für alle definierten I/O-Testfälle dieselben Ausgaben wie die Quellimplementierung? Dieses Kriterium ist formal als I/O-Äquivalenz definiert: Für identische Eingaben müssen Quell- und Zielimplementierung identische Ausgaben produzieren. Eniser et al.\ entwickeln mit Fluorine einen cross-language Differenzfuzzer, der diese Äquivalenz automatisch prüft; für den vorliegenden PoC erfolgt der Abgleich manuell an repräsentativen Testfällen \parencite{eniser2025realworld}.

\textbf{Vollständigkeit (Feature Parity):} Sind alle acht identifizierten Funktionalitäten im Zielcode implementiert? Nicht abgedeckte Funktionalitäten gelten als Parity-Lücken.

\textbf{Wartbarkeit:} Ist der generierte Code lesbar, strukturiert und dokumentiert? Als Maßstab dienen messbare Wartbarkeitskriterien wie Methodenlänge, zyklomatische Komplexität und Kommentartiefe \parencite{nunes2025maintainability}. Code, der zwar korrekt, aber kaum wartbar ist, gilt als Qualitätsmangel.

\textbf{Idiomatizität:} Folgt der Code den Konventionen des Zielstacks? Für Go bedeutet das: explizite Fehlerbehandlung, stdlib-first, kurze Receiver-Namen, Interfaces an der Consumer-Seite. Der Go-Craftsman-Skill operationalisiert diese Kriterien als implementierungsbegleitende Checkliste. Dass LLMs auch bei korrekt übersetztem Code in 1--15 Prozent der Fälle nicht-idiomatischen Stil produzieren, belegen Eniser et al.\ mit einer Clippy-Analyse generierter Rust-Übersetzungen \parencite{eniser2025realworld}; ein analoges Muster ist für Go-Generierungen plausibel und wird im Evaluationsdurchlauf mit \texttt{go vet} und manuellem Code-Review geprüft.

\textbf{Rückverfolgbarkeit zur Spec:} Lässt sich jeder implementierte Codeabschnitt auf ein Akzeptanzszenario in \texttt{SPEC.md} zurückführen? Die Inline-Kommentare im generierten Code (z.\,B. \texttt{// FR-001/FR-002}) ermöglichen diese Prüfung direkt im Code-Review.

\subsection{Sicherstellung der Vollständigkeit}

Ein strukturelles Risiko agentischer Pipelines ist das vorzeitige Beenden eines Laufs, bevor alle Funktionalitäten tatsächlich verarbeitet wurden. Dies adressiert das in Kapitel~\ref{implementierung} beschriebene Completion-Promise-Prinzip: Ein subjektives Fertigkeitsgefühl des Orchestrators gilt erst dann als bestätigt, wenn die Verifier-Prüfung alle Artefakte auf Disk verifiziert hat. Zur Sicherstellung der Vollständigkeit werden im PoC drei komplementäre Mechanismen kombiniert.

Erstens dient \texttt{FUNCTIONALITY-INDEX.md} als verbindliche Vollständigkeitsliste für Phase~1: Jedes identifizierte Funktionalitätselement trägt einen expliziten Status (\emph{ready}, \emph{needs-clarification}, \emph{blocked}), und der Orchestrator darf Phase~2 erst beginnen, wenn für jedes \emph{ready}-Element eine \texttt{SPEC.md} im zugehörigen Unterordner vorliegt. Der Orchestrator-Completion-Skill operationalisiert diese Prüfung als artefaktbasierte Checkliste, die vor jedem Phasenübergang durchlaufen werden muss.

Zweitens fungiert \texttt{TASKS.md} als Vollständigkeitsliste für Phase~3: Jede Aufgabe ist als Checkbox modelliert, und der Builder hakt sie erst ab, wenn der zugehörige Code und ein Verifikationsnachweis in \texttt{BUILD.md} vorliegen. Eine \texttt{TASKS.md} mit offenen Checkboxen signalisiert dem Orchestrator, dass Phase~3 noch nicht abgeschlossen ist.

Drittens schreibt der Spec-Quality-Skill vor, dass eine \texttt{SPEC.md} keine stillen Annahmen enthalten darf: Alle unklaren Punkte müssen explizit als offene Fragen markiert sein. Damit ist auch auf Spezifikationsebene keine implizite Vollständigkeit möglich, die den Evaluationsschritt verfälschen könnte.

Die Ergebnisse der Evaluation fließen in die abschließende Bewertung und den Ausblick in Kapitel~\ref{fazit} ein.
